\documentclass[10pt]{article}

\usepackage[utf8]{inputenc}

\usepackage[T1]{fontenc}

\usepackage{amsmath}

\usepackage{amsfonts}

\usepackage{amssymb}

\usepackage[version=4]{mhchem}

\usepackage{stmaryrd}

\usepackage{bbold}

\usepackage{graphicx}

\usepackage[export]{adjustbox}

\graphicspath{ {./images/} }



\begin{document}

равносильны, при этом $\left|f^{*}(\xi)\right|=\left|F^{*}(\xi)\right|=\chi(\xi)$ почти всюду на $T$. Внутренние функции $G(z)$, имеющие вит $G(z)=B(z) \cdot S(z)$, полностью характеризуются условиями $|G(z)|<1$ в $D$ и $\left|G^{*}(\xi)\right|=1$ почти всюду на $T$. Часто используют разложение $f(z)=\sqrt{f_{0}(z)} \cdot\left(\sqrt{f_{0}(z)} B(z)\right)$ произвольной функции $f(z) \in H^{1}$ в произведение двух функций из $H^{2}$ (см. [4], [5]).



Класс $H^{2}$ занимает особое место среди X. к., так как является гильбертовым пространством с воспроизводящим ядром и имеет простое ошисание через коэф্фициенты Теӥлора:



\[

\sum_{n=0}^{\infty} a_{n} z^{n} \in H^{2} \Leftrightarrow\left\{a_{n}\right\} \in l^{2} .

\]



Важную роль сыграло изучение оператора умножения на $z$, или о пе р а т о р а с д в и г а, в піространстве $H^{2}$; оказалось, что все инвариантные подпространства этого олератора порождены внутренними функциями $G(z)$, т. с. Имеют вид $G(z) \cdot H^{2}$ (cM. [4]).



Относительно поточечного умножения и sup-нормы класс $H^{\infty}$ является банаховой алгеброй с весьма сложным строснием пространства максимальных идеалов $\mathfrak{M}$ и границы Шилова (см. [4]); вопрос о плотности идеалов $(z-\xi) \cdot H^{\infty}, \xi \in D$, в пространстве $\mathfrak{M}$ с обычной топологисй Гельфанда (т. н. ін о о л с ма ко роны) был релчен положительно на основе описания униве рсальных интер поляционных послед о в а т е ль но с т е й - госледовательностей $\left\{z_{n}\right\}$, $z_{n} \in D$, таких, что $\left\{\left\{f\left(z_{n}\right)\right\}: f(z) \in H^{\infty}\right\}=l^{\infty}$ (см. [5]).



Х. к. $H^{p}(G), p>0$, алалитич. функций $f(z)$ в областях $G \subset \mathbb{C}$, отличных от круга, можно определить (в общем случае неэквивалентно) исходя либо из условия существования у функций $|f(z)|^{p}$ гармонич. мажоранты в $G$, либо из условия ограниченности интегралов $\int L_{r}|f(z)| p|d z|$ по семействам контуров $\left\{L_{r}\right\}$, $L_{r} \subset G$, в каком-то смысле приближающих границу области $G$. Первый способ позволяет определить также X. к. на римановых поверхностях. Второй способ̄ приводит к классам, лучше приспособленным для решения әкстремальных и аппроксимационных задач; в случае жордановых областей $G$ со спрямляемой границей последние классы наз. к л а с с а м и С м и р н о в а и обозначаются $E^{p}(G)$ (см. [2]). Для полуплоскости, напр. $\boldsymbol{P}=\{\operatorname{Re} z>0\}$, глассы $H^{p}(\boldsymbol{P}), \quad p>0$, огределяемые условием



\[

\sup _{x>0} \int_{-\infty}^{\infty}|f(x+i y)|^{p} d y<\infty

\]



по свойствам близки к Х. к. для круга, однако их приложения в гармонич. анализе связаны уже не с теорией рядов Фурье, а с теорией преобразований Фурье.



Х. к. аналитич. функций $f(z)=f\left(z_{1}, \ldots, z_{n}\right)$ в единичном ппаре $B^{n}$ и единичном поликруге $D^{n}$ пространства $\mathbb{C}^{n}$ огределяются условием (;) с заменой окружности $T$ соответственно сферой $\partial B^{n}$ или остовом $T^{n}$ полікруга. Специфика многомерного случая проявляется прежде всего в отсутствии простой характеристики множеств нулей и факторизации функций соответствующих Х. к. (см. [6], [10]). Х. к. определяются, причем различными способами, и для других областей в $\mathbb{C}^{n}$ (см. [10]).



Многомерными аналогами X.к. (см. [3]) являются т. н. П р о с т а н с т в а $\mathrm{X}$ а р ди- пространства $H^{p}\left(\mathbb{R}^{n}\right), \quad p>0, \quad$ систем Рисса - действительных вектор-функций $F(x, y)=\left\{u_{j}(x, y)\right\}_{j=0}^{n}, \quad x \in \mathbb{R}^{n}, \quad y>0$, удовлетворяющих обобщен ным условия я $\mathrm{K}$ оIII и-Р и м а на



\[

\sum_{j=0}^{n} \frac{\partial u_{j}}{\partial x_{j}}=0, \quad \frac{\partial u_{j}}{\partial x_{k}}=\frac{\partial u_{k}}{\partial x_{j}}, \quad x_{0}=y

\]



\begin{center}

\includegraphics[max width=\textwidth]{2023_07_09_808c5c62fb98326cff59g-1}

\end{center}



Определение этих пространств можно дать и в терминах лишь «действительных частей» $u_{0}(x, y)$ систем $F(x, y)$, потребовав, ттобы функция $u_{0}(x, y)$ была гармонической, a еe максимальная фуунция по некасатольным путям



$M_{a} u(x)=\sup _{t, y}\left\{\left|u_{0}(t, y)\right|:|t-x|<a y\right\} \in L^{p}\left(\mathbb{R}^{n}\right), a>0$.



Јри $p>1$ переход от функции $u_{0}(x, y)$ к ее граничным значениям дает отождествление пространств $H^{p}\left(\mathbb{R}^{n}\right)$ и $L^{p}\left(\mathbb{R}^{n}\right)$, поэтому интерес представляет лишь случай $p \leqslant 1$. Именно в рамках пространств $H^{p}\left(\mathbb{R}^{n}\right)$ были первоначально установлены такие фундаментальные результаты тоории X. к., как реализация сопряженного пространства $\left(H^{1}\right)^{*}$ в виде Ігространства функций ограниченного среднего колебания (см. [8], [9]) и атомич. разложение классов $H^{p}, p \leqslant 1$ (см. [7]). Характеристика Х. к. в терминах максимальной функции в ряде случаев требует привлечения вероятностных понятий, связанных с броуновским движением (см. [8]).



Абстрактные X.к. возникают в теории равномерных алгебр и не связаны непосредственно с аналитич. функциями. Фиксируется замкнутая алгебра $A$ непрерывных фушкций на компакте $X$ и нек-рый гомоморфизм $\varphi$ : $A \longrightarrow \mathbb{C} ;$ существует положительная мера $d \mu$ на $X$, представляющая $\varphi: \varphi(f)=\int_{X} f \cdot d \mu, f \in A$. По определению, классы $H^{p}(d \mu), 0<p \leqslant \infty$, суть замыкания (слабое замыкание при $p=\infty$ ) алгебры $A$ в пространствах $L^{p}(d \mu)$; изучение классов $H^{p}(d \mu)$ повволяет получить дополнительную информацию об алгебре $A$ (cм. [11]).



Jum.: [1] R i e s z F., "Math. Z.», 1923, Bd 18, S. 87-95; [2] П р и в а л о в и. И., Граничные свойства аналитических Функций, 2 изд., M.- JI., 1950: [3] S t е i n E., W e i s s G.,

"Acta math.", 1960 , v. 103, p. 25-62; [4] Г о ф м а н K., Банахо"Acta math.", 1960, v. 103, p. $25-62 ;[4]$ Г оф м а н К., Банахо[5] D u r e n P., Theory of $H^{p}$ spaces, N. Y.- L., 1970; [6] P yд и пІ У., Теория функций в поликруге, пер. с англ., М., 1974; 77 C o i f m a R. R., W e is s G., "Bull. Amer. Math. Soc.»,

1977, v. 83, No 4, p. 569-645; [8] P e t e r e n K. E., Brownian motion, Hardy spaces and bounded mean oscillation, L., 1977; [9] K o o s is P., Introduction to $H_{p}$ spaces. With an appendix on Wolff's proof of the corona theorem, Camb., 1980; [10] R u d i n W., Function theory in the unit ball of $C^{n}$, ІІер. с англ., M., 1973 .



С. В. Шведенко.



ХАРДИ НЕРАВЕНСТВО д л я р я д о в: если $p>1$, $a_{n} \geqslant 0$ и $A_{n}=a_{1}+\ldots+a_{n}, n=1,2, \ldots$, то



\[

\sum_{n=1}^{\infty}\left(\frac{A_{n}}{n}\right)^{p}<\left(\frac{p}{p-1}\right)^{p} \sum_{n=1}^{\infty} a_{n}^{p}

\]



кроме случая, когда все $a_{n}$ равны нулю. Константа $\left(\frac{p}{p-1}\right)^{p}$ в этом неравенстве наилутшая.



Х. н. для интегралов



$\int_{0}^{+\infty} x-p\left|\int_{0}^{x} f(t) d t\right|^{p} d x<\left(\frac{p}{p-1}\right)^{p} \int_{0}^{+\infty}|f(x)|^{p} d x, \quad p>1$



$\mathbf{n}$



$\int_{0}^{+\infty}\left|\int_{x}^{+\infty} f(t) d t\right|^{p} d x<p p \int_{0}^{+\infty}|x f(x)|^{p} d x, \quad p>1$.



Неравенства сграведливы для всех фунгций, для к-рых конечны правые части неравенств, кроме случая, когда функция $f$ почти всюду на интервале $(0,+\infty)$ равна нуліо (в этом случае неравенства обращаются в равен-





\end{document}